% !TeX root = ../sections/03_solutions.tex
\subsection{flow}

\begin{background}
	この試験では，次の三つの力が中心になっている。
	
	\begin{enumerate}
		\item 周期関数を分類する力
		\item フーリエ係数を計算する力
		\item 内積空間の考え方で近似・直交化を扱う力
	\end{enumerate}
	
	特に，フーリエ級数と平均 \(2\) 乗誤差は別々の話ではない。
	フーリエ級数は，三角関数の張る空間への直交射影として理解できる。
	そのため，平均 \(2\) 乗誤差を最小にする近似は，
	フーリエ係数によって与えられる。
\end{background}

\begin{strategy}
	問題を見たときは，次の順番で判断する。
	
	\begin{enumerate}
		\item 周期関数の分類問題では，端点と特異点を見る。
		\item フーリエ級数では，まず偶関数か奇関数かを見る。
		\item 偶関数なら余弦項だけ，奇関数なら正弦項だけを考える。
		\item 区分的に定義された関数では，積分区間を分割する。
		\item gnuplot では，関数・係数・部分和を順番に定義する。
		\item 平均 \(2\) 乗誤差の最小化では，フーリエ係数を使う。
		\item シュミットの直交化法では，内積を使って前の成分を引く。
	\end{enumerate}
	
	この試験の対策としては，まず次の公式を確実に覚える。
	
	\[
	a_0=\frac{1}{\pi}\int_{-\pi}^{\pi}f(x)\,dx,
	\]
	\[
	a_n=\frac{1}{\pi}\int_{-\pi}^{\pi}f(x)\cos nx\,dx,
	\]
	\[
	b_n=\frac{1}{\pi}\int_{-\pi}^{\pi}f(x)\sin nx\,dx.
	\]
	
	次に，内積の公式
	\[
	(f,g)=\int_a^b f(x)g(x)\,dx
	\]
	と，シュミットの直交化法
	\[
	\varphi_k
	=
	\psi_k
	-
	\sum_{j=1}^{k-1}
	\frac{(\psi_k,\varphi_j)}{(\varphi_j,\varphi_j)}
	\varphi_j
	\]
	を覚える。
	
	この二つを結びつけて理解できると，
	フーリエ級数，平均 \(2\) 乗誤差，直交関数の問題が同じ構造で見える。
\end{strategy}